\section{The Coulomb system and its small-coupling limit}
\label{sec:setup}

Let
\[
 \mathfrak h=L^2(\R^3;\C^2),\qquad
 \mathfrak H_m=\bigwedge\nolimits^m\mathfrak h
\]
be the one-electron and antisymmetric \(m\)-electron spaces.  Spin is
expressed in the basis \(\uparrow,\downarrow\), and spatial reflections
and rotations act trivially on spin.  For \(D>0\), put
\begin{equation}
 u_D(x)=-\frac1{|x-De_z|}-\frac1{|x+De_z|},
 \qquad h_D=-\frac12\Delta+u_D,\qquad e_z=(0,0,1).
 \label{eq:unit-charge-potential}
\end{equation}
The corresponding \(m\)-electron operators are
\begin{equation}
 H_{m,\lambda}
 =\sum_{i=1}^m h_D(x_i)
   +\lambda\sum_{1\le i<j\le m}\frac1{|x_i-x_j|},
 \qquad m=2,3,\quad \lambda\ge0 .
 \label{eq:scaled-coulomb-hamiltonian}
\end{equation}
They act in \(\mathfrak H_m\).  Hardy's inequality in a particle
coordinate, or in a relative coordinate, gives
\[
 \bigl\||x_i-R|^{-1}\psi\bigr\|_2+
 \bigl\||x_i-x_j|^{-1}\psi\bigr\|_2
 \le C\|\nabla\psi\|_2
 \le \varepsilon\|\Delta\psi\|_2+C_\varepsilon\|\psi\|_2 .
\]
Thus these operators are self-adjoint on the antisymmetric part of
\(H^2(\R^{3m};\C^{2^m})\), and their form domain is the corresponding
\(H^1\) space.  The same estimates show that \(\lambda\mapsto
H_{m,\lambda}\) is an analytic family with this common operator domain.
We write \(E_m(\lambda)\) for its ground energy whenever the ground
level under consideration exists.

For a normalized three-electron state, its spin-summed density is
\begin{equation}
 n_\Psi(x)
 =3\sum_{\sigma_1,\sigma_2,\sigma_3}
   \int_{\R^6}
   |\Psi(x,\sigma_1,y_1,\sigma_2,y_2,\sigma_3)|^2\,dy_1\,dy_2 .
 \label{eq:density-definition}
\end{equation}
In particular, \(\int n_\Psi=3\).  Ensemble densities are defined by
linearity, as detailed in \Cref{sec:ensemble}.

\Needspace{5\baselineskip}
\begin{proposition}[Dilation to integer nuclear charges]
\label{prop:scaling}
Suppose that the density of a ground state of \(H_{3,1/Z}\), for a
positive integer \(Z\), has no noninteracting ground-state ensemble
representation with a real spin-independent potential in
\(L^{3/2}(\R^3)+L^\infty(\R^3)\).
The physical Coulomb system with three electrons and charges \(Z,Z\)
at
\[
 (0,0,D/Z),\qquad (0,0,-D/Z)
\]
then has a ground-state density with the same obstruction.
\end{proposition}

\begin{proof}
Use physical coordinates \(\mathbf r_i\) and set \(x_i=Z\mathbf r_i\).
The unitary dilation
\[
 (\mathcal U_Z\psi)(x_1,\ldots,x_m)
 =Z^{-3m/2}\psi(x_1/Z,\ldots,x_m/Z)
\]
leaves spin unchanged.  It transforms the physical Coulomb Hamiltonian
into \(Z^2H_{m,1/Z}\).  In particular, normalizability, the ground-space
property, and antisymmetry are preserved, and
\[
 n_{\rm phys}(\mathbf r)=Z^3 n_{\rm scaled}(Z\mathbf r).
\]
A physical noninteracting representing potential would transform into
\[
 \widetilde v(x)=Z^{-2}v(x/Z).
\]
For a decomposition \(v=v_1+v_\infty\), with
\(v_1\in L^{3/2}\) and \(v_\infty\in L^\infty\), the corresponding
pieces satisfy
\[
 \|Z^{-2}v_1(\,\cdot\,/Z)\|_{3/2}=\|v_1\|_{3/2},
 \qquad
 \|Z^{-2}v_\infty(\,\cdot\,/Z)\|_\infty
 =Z^{-2}\|v_\infty\|_\infty .
\]
The dilation of the representing density matrix preserves its trace
and support in the ground eigenspace.  It would therefore give the
excluded representation of the scaled density.
\end{proof}

\subsection{The two lowest spatial orbitals}

All subsequent one-particle orbitals are spatial orbitals.  For fixed
\(D\) we abbreviate \(u_D\) to \(u\).  Reflection symmetry will mean
reflection in \(z=0\), with \(x=(x_1,x_2,z)\).

\begin{proposition}[Two-centre spectrum]
\label{prop:two-centre-spectrum}
For \(D\ge100\), the two lowest eigenvalues of \(h_D\) are simple.
They have normalized real, axisymmetric eigenfunctions \(b,e\), where
\(b>0\) is even and \(e\) is odd and strictly positive on \(z>0\).
Write
\[
 \epsilon_b=-\frac{\gamma^2}{2},\qquad
 \epsilon_e=-\frac{k_0^2}{2},\qquad
 d_*=\epsilon_e-\epsilon_b .
\]
Then
\begin{equation}
 0<d_*<\frac{100}{D},\qquad
 |\gamma-1|+|k_0-1|<\frac{200}{D}.
 \label{eq:two-centre-numerics}
\end{equation}
The spectrum on \(\{b,e\}^{\perp}\) lies more than \(0.2\) above
\(\epsilon_e\).  In particular, the spectral gap above \(b\) in the
even subspace is greater than \(0.2\).
\end{proposition}

\begin{proof}
Let \(\varphi(x)=\pi^{-1/2}e^{-|x|}\) be the hydrogenic ground state.
We recall the two hydrogenic bounds needed for localization.  After
spherical-harmonic decomposition, the radial operators and first-order
factors are
\[
 H_\ell=-\frac12\frac{d^2}{d\rho^2}
       +\frac{\ell(\ell+1)}{2\rho^2}-\frac1\rho,
 \qquad
 A_\ell=\frac d{d\rho}-\frac{\ell+1}{\rho}+\frac1{\ell+1}.
\]
Here \(H_\ell\) has its Friedrichs realization in
\(L^2((0,\infty),d\rho)\), and \(A_\ell\) is closed from
\(C_c^\infty(0,\infty)\).
On the corresponding Friedrichs form domains, integration by parts gives
\[
 \frac12A_\ell^*A_\ell=H_\ell+\frac1{2(\ell+1)^2},
 \qquad
 \frac12A_\ell A_\ell^*=H_{\ell+1}+\frac1{2(\ell+1)^2}.
\]
The first identity and centrifugal comparison give
\(H_\ell\ge-1/8\) for \(\ell\ge1\).  Also
\(\ker A_0=\operatorname{span}\{\rho e^{-\rho}\}\), and the second
identity gives \(A_0A_0^*\ge3/4\).  The polar decomposition transfers
this last bound to \(A_0^*A_0\) on \((\ker A_0)^\perp\).
Thus the hydrogenic operator has ground energy \(-1/2\) and is
bounded below by \(-1/8\) on \(\varphi^\perp\); the radial
\(\ell=1\) function \(\rho^2e^{-\rho/2}\) attains \(-1/8\).

Choose a Lipschitz angle \(\tau(\rho)\) that is zero for
\(\rho\le D/4\), is \(\pi/2\) for \(\rho\ge D/2\), and is linear
between these radii.  Its slope is \(2\pi/D\).  Around the two nuclei
use cutoffs
\(\chi_\pm(x)=\cos\tau(|x\mp De_z|)\), set equal to zero outside their
respective balls of radius \(D/2\), and complete them by the
nonnegative cutoff \(\chi_{\rm out}\) to a quadratic partition:
\[
 \chi_+^2+\chi_-^2+\chi_{\rm out}^2=1.
\]
The two balls are disjoint.  On either transition annulus the
nonconstant pair is a cosine and a sine, so
\begin{equation}
 \frac12\sum_{\nu\in\{+,-,{\rm out}\}}|\nabla\chi_\nu|^2
 \le \frac{2\pi^2}{D^2}<\frac{20}{D^2}.
 \label{eq:two-centre-ims-cost}
\end{equation}
Lipschitz cutoffs suffice for the form identity; alternatively one may
smooth the ramp with arbitrarily small increases in these constants.

On either inner ball, the attraction of the other nucleus has
magnitude at most \(2/(3D)<1/D\).  On the support of
\(\chi_{\rm out}\), both nuclear distances are at least \(D/4\);
there the total attraction is at most \(8/D\).  Applying the IMS
identity and the hydrogenic ground-energy bound gives
\begin{equation}
 h_D\ge-\frac12-\frac1D-\frac{20}{D^2}
 \quad\text{as a quadratic form}.
 \label{eq:two-centre-bottom-lower}
\end{equation}
Here and below the exterior lower bound \(-8/D\) is larger than the
displayed inner lower bound for \(D\ge100\).

If a test function \(\psi\) is orthogonal to both
\(\chi_\pm\varphi(\,\cdot\,\mp De_z)\), then
\(\chi_\pm\psi\) is orthogonal to the ground state of the corresponding
hydrogenic operator.  The next hydrogenic energy is \(-1/8\).
Consequently,
\begin{equation}
 \langle\psi,h_D\psi\rangle
 \ge\left(-\frac18-\frac1D-\frac{20}{D^2}\right)\|\psi\|_2^2
 \label{eq:two-centre-third-lower}
\end{equation}
on a subspace of codimension two.

For the upper bound, use the two disjointly supported trial functions
\(\chi_\pm\varphi(\,\cdot\,\mp De_z)\).  Their squared norms exceed
\(1/2\), since the hydrogenic mass outside radius \(D/4\) is
\[
 e^{-D/2}\left(1+\frac D2+\frac{D^2}{8}\right)<\frac12.
\]
The ground-state identity for the hydrogenic operator and the bound
\(|\nabla\chi_\pm|\le2\pi/D\) show that their Rayleigh quotients for
\(h_D\) are at most
\begin{equation}
 -\frac12+\frac{4\pi^2}{D^2}
 <-\frac12+\frac{40}{D^2}.
 \label{eq:two-centre-two-upper}
\end{equation}
Their even and odd combinations give the same bound in the two
parity subspaces.  The Coulomb potential is relatively compact with
respect to the free resolvent: one first truncates at infinity and
near the nuclei, and uses local compactness for the bounded remainder.
The omitted nuclear piece has operator norm from \(H^2\) to \(L^2\)
of order \(\sqrt{\varepsilon}\) when cut off to a ball of radius
\(\varepsilon\), by \(H^2\hookrightarrow L^\infty\) and
\(\||x|^{-1}\mathbf 1_{|x|<\varepsilon}\|_2=O(\sqrt\varepsilon)\).
The far-field remainder tends to zero in \(L^\infty\).
Hence the essential spectrum is \([0,\infty)\).
The min--max principle now gives two negative levels, one in each
parity, and no third level below the lower bound in
\Cref{eq:two-centre-third-lower}.

A full-space ground state can be chosen nonnegative by the
absolute-value inequality for its kinetic energy.  It is strictly
positive, including at the nuclei.  One direct way to include those
points is to use
\[
 b=\left(-\frac12\Delta+c\right)^{-1}
       (c+\epsilon_b-u)b,\qquad c>|\epsilon_b|,
\]
whose resolvent kernel is strictly positive; the right-hand side is
nonnegative and nonzero, and \(b\in H^2(\R^3)\) is continuous.
The same argument with the Dirichlet resolvent on \(z>0\) gives
strict positivity of its ground state in that half-space.
Positivity and the usual ground-state simplicity argument identify
the low even level with \(b\) and the low odd level with the odd
extension of that Dirichlet ground state.  Simplicity also follows
from the two-level count above.  The full-space ground state cannot
be odd, so \(\epsilon_b<\epsilon_e\).  Rotation about the \(z\)-axis
preserves the unique normalized positive full-space and half-space
ground states, proving axial symmetry.

Combining \Cref{eq:two-centre-bottom-lower,eq:two-centre-two-upper}
gives
\[
 d_*\le\frac1D+\frac{60}{D^2}<\frac{100}{D}.
\]
The same two bounds place \(\gamma\) and \(k_0\) within
\(200/D\) of \(1\), with the stated bound on their sum.  Finally,
\Cref{eq:two-centre-third-lower,eq:two-centre-two-upper} give the
gap estimate
\[
 \frac38-\frac1D-\frac{60}{D^2}>0.2 .
\]
Since \(e\) is odd, the same lower bound holds above \(b\) in the
even subspace.
\end{proof}

\subsection{Ground states and perturbation}

\begin{proposition}[Small repulsive coupling]
\label{prop:small-coupling}
Fix \(D\ge100\).  For all sufficiently small \(\lambda\ge0\), the
ionic operator \(H_{2,\lambda}\) has a simple isolated ground state,
and \(H_{3,\lambda}\) has a simple isolated ground state
\(\Psi_\lambda\) in the sector \(S_z=1/2\).  Its energy is an absolute
ground energy of \(H_{3,\lambda}\).

The ionic ground state is a spin singlet, is even under simultaneous
reflection of its two spatial coordinates, and is axisymmetric.
The state \(\Psi_\lambda\) is real, axisymmetric, and odd under
simultaneous reflection of its three spatial coordinates.
The states and energies can be chosen analytic near zero; in
particular, their state vectors converge with error \(O(\lambda)\)
in \(H^2\).  At zero coupling,
\begin{equation}
 \Psi_0=b\uparrow\wedge b\downarrow\wedge e\uparrow,\qquad
 E_2(0)=2\epsilon_b,\qquad
 E_3(0)=2\epsilon_b+\epsilon_e .
 \label{eq:zero-coupling-states}
\end{equation}
The ionic state has spatial factor \(b(x_1)b(x_2)\) and the normalized
singlet spin factor.
\end{proposition}

\begin{proof}
At \(\lambda=0\), the Pauli principle fills the spin orbitals in
\Cref{eq:zero-coupling-states}.  In the sector \(S_z=1/2\), moving
the doubly occupied spatial orbital from \(b\) to \(e\) costs
\(d_*>0\); using any further spatial orbital costs a positive gap
by \Cref{prop:two-centre-spectrum}.  The displayed determinant is
therefore a simple isolated ground state in that sector.  The ionic
ground state is likewise simple and isolated.  The noninteracting
ionization thresholds are \(2\epsilon_b\) and \(\epsilon_b\),
respectively, and are strictly above these ground energies.

The common-domain analytic perturbation argument applies to each
isolated rank-one spectral projection; see
\cite[Chapter~VII, \S1.3, Theorem~1.7, and \S2.2, Theorem~2.6]{Kato}.
To include the graph-norm conclusion explicitly, write
\(V_{\rm ee}=\sum_{i<j}|x_i-x_j|^{-1}\) and
\(R_0(\zeta)=(H_{m,0}-\zeta)^{-1}\).  On a contour enclosing only
the relevant zero-coupling level,
\[
 (H_{m,\lambda}-\zeta)^{-1}
 =R_0(\zeta)\bigl(I+\lambda V_{\rm ee}R_0(\zeta)\bigr)^{-1}.
\]
Here \(R_0(\zeta):L^2\to H^2\) and
\(V_{\rm ee}R_0(\zeta):L^2\to L^2\) are bounded uniformly on the
contour.  A Neumann series therefore makes the resolvent analytic
as a map into \(H^2\).  Integrating around the contour, applying the
resulting projection to the unperturbed state, and normalizing gives
a real-analytic state in that space.  Choose its
real phase by positive overlap with the zero-coupling state.
Spectral separation, or the min--max principle with the nonnegative
repulsion, keeps this branch the lowest level of its sector.

It remains to check the absolute three-electron ground property.
In the sectors \(S_z=\pm3/2\), all three spatial orbitals must be
distinct.  Their zero-coupling energy is at least
\(\epsilon_b+\epsilon_e+\epsilon_3\), where
\(\epsilon_3=\inf\sigma(h_D|_{\{b,e\}^\perp})\); it exceeds
\(2\epsilon_b+\epsilon_e\) by more than \(0.2\).
Nonnegative repulsion preserves this lower bound, whereas the
continued \(S_z=1/2\) energy converges to
\(2\epsilon_b+\epsilon_e\).  Spin flip identifies the two
\(S_z=\pm1/2\) ground energies.  These four sectors exhaust the
three-electron spin space, proving the claim.

All the stated symmetries commute with the Hamiltonian and its
spectral projection.  Their eigenvalues at zero coupling persist on
the simple branch.  For axial rotations, one may equivalently note
that the projection of the invariant zero-coupling state is still
invariant.  The ionic singlet subspace is also invariant, so its
spin character persists.
\end{proof}

For the rest of the small-coupling argument define
\begin{equation}
 k=k_\lambda=\sqrt{2\bigl(E_2(\lambda)-E_3(\lambda)\bigr)},\qquad
 A(r)=e^{-kr}r^\beta,\qquad
 \beta=\frac{2-2\lambda}{k}-1 .
 \label{eq:ionization-parameters}
\end{equation}
Thus \(k_\lambda\to k_0\).  Once \(D\) is fixed, the perturbation
interval may be shortened so that
\begin{equation}
 \gamma-k_\lambda\ge\frac12(\gamma-k_0)>0 .
 \label{eq:uniform-fast-slow-separation}
\end{equation}
No lower bound uniform in \(D\) on the length of this interval is
asserted or needed.

\begin{lemma}[Density regularity and compact convergence]
\label{lem:density-regularity}
Regard
\(F_{\rm true}(x)=\sqrt3\,\Psi_\lambda(x,\cdot)\) as a
Hilbert-space-valued function, retaining the removed-electron spin
and all ionic coordinates and spins.  Then
\[
 F_{\rm true}\in H^2(\R^3;\mathcal K),\qquad
 n_\lambda=\|F_{\rm true}\|_{\mathcal K}^2
       \in H^2(\R^3)\cap L^1(\R^3),\qquad
 \sqrt{n_\lambda}\in H^1(\R^3),
\]
where \(\mathcal K\) is the corresponding fiber Hilbert space.
Moreover \(\Delta n_\lambda\in L^1(\R^3)\).
These assertions hold for every \(H^2\) three-electron eigenstate,
with no small-coupling assumption.

For the branch in \Cref{prop:small-coupling},
\(n_\lambda\to n_0=2b^2+e^2\) in \(H^2\) and \(L^1\), hence
uniformly on compact sets; the densities are commonly tight as
\(\lambda\downarrow0\).  Each fixed compact set has a positive
lower bound for \(n_\lambda\) for sufficiently small \(\lambda\).
The density is even and axisymmetric.
\end{lemma}

\begin{proof}
Fubini's theorem applied to the derivatives in the distinguished
coordinate gives the Hilbert-valued \(H^2\) assertion.  In three
dimensions the corresponding Sobolev estimates yield
\[
 \|F_{\rm true}\|_{L^\infty_x(\mathcal K)}
 +\|\nabla_xF_{\rm true}\|_{L^4_x(\mathcal K)}
 \le C\|F_{\rm true}\|_{H^2_x(\mathcal K)} .
\]
The distributional product identities
\begin{align}
 \partial_j n_\lambda
 &=2\operatorname{Re}
      \langle F_{\rm true},\partial_jF_{\rm true}\rangle_{\mathcal K},
 \notag\\
 \partial_i\partial_j n_\lambda
 &=2\operatorname{Re}
      \langle\partial_iF_{\rm true},\partial_jF_{\rm true}\rangle_{\mathcal K}
   +2\operatorname{Re}
      \langle F_{\rm true},\partial_i\partial_jF_{\rm true}\rangle_{\mathcal K}
 \label{eq:density-product-identities}
\end{align}
then prove \(n_\lambda\in H^2\).  The same identities and
Cauchy--Schwarz, now using \(L^2_x\) on both factors, give
\(\Delta n_\lambda\in L^1\).  The norm chain rule gives
\begin{equation}
 \int_{\R^3}|\nabla\sqrt{n_\lambda}|^2
 \le \int_{\R^3}\|\nabla_xF_{\rm true}\|_{\mathcal K}^2 .
 \label{eq:density-kinetic-bound}
\end{equation}

The product estimates also show continuity of
\(F\mapsto\|F\|_{\mathcal K}^2\) from \(H^2_x(\mathcal K)\) to
\(H^2_x\).  Thus the graph-norm convergence in
\Cref{prop:small-coupling} gives the claimed density convergence.
For the \(L^1\) convergence, use
\[
 \|n_\lambda-n_0\|_1
 \le
 \bigl(\|F_{{\rm true},\lambda}\|_2+\|F_{{\rm true},0}\|_2\bigr)
 \|F_{{\rm true},\lambda}-F_{{\rm true},0}\|_2 .
\]
It implies common tightness.  Since \(b\) is continuous and
strictly positive everywhere, \(n_0\ge2b^2\) has a positive minimum
on each compact set.  Uniform convergence gives the corresponding
compact lower bound for \(n_\lambda\).  The spatial symmetries
follow by taking the squared fiber norm.
\end{proof}

The last lemma makes no claim of global positivity from compact
convergence alone.  The exterior argument below supplies positivity
outside a fixed ball uniformly for sufficiently small coupling, and
thereby supplies the global positivity needed for the angle
construction.
